\section{Introduction}
\label{sec:introduction}

If a restricted quantum record admits a common effective update on every transition observed so far, what has actually been learned about its future dynamics?

The existence of a physical map fitting a finite dataset is an important consistency statement, but it is not yet a predictive law or a point-identification result.  Finite observations can leave a non-singleton feasible model set, and representatives of that set can differ on directions or transitions not constrained by the observations.  Conversely, failure of one selected representative does not by itself show that the model class is misspecified.  These distinctions are standard in set-membership system identification, where finite data and uncertainty assumptions define a feasible set before any representative model is selected for prediction \cite{MilaneseNovara2004,MilaneseNovara2011}.  They also arise directly in incomplete quantum process tomography: incomplete process information can leave multiple channels compatible with the data, after which an additional selection principle may be used to choose a single estimator \cite{Ziman2008,TeoEtAl2011}.

This paper isolates the same logical distinction in a finite quantum coarse-dynamics setting where compatibility, selection of one witness, and later prospective testing can be certified separately.  We work with a fixed $N=8$ quantum system and a fixed proper coarse record denoted $R2_{07}$.  The admissible model class is deliberately broad: one physical, source-independent, time-independent, whole-algebra CPTP map on the full record algebra, with all central-sector transfers allowed.  The discrepancy on each transition is the half diamond norm of the difference between the next complete source-to-record channel and the channel obtained by applying the common map to the current source-to-record channel.

The reviewed finite sequence is as follows.  First, four already exposed transitions admit an explicit common map $\Phi_4$ with a rigorous error certificate below $10^{-2}$.  Second, that exact map is frozen, and two later transitions are chosen without access to their physical outcomes by an outcome-blind schedule plus an independent historical-exposure audit.  On those pristine transitions the frozen map fails by a large margin: independently reviewed lower bounds exceed $0.62$ on each transition.  Third, those failed predictions are added to the fitting constraints and the \emph{same full physical map class} is considered on the enlarged six-transition dataset.  A different explicit common map $\Phi_6$ fits all six transitions, and hostile review gives the rigorous minimax bracket
\begin{equation}
0.002306082156934015\ldots
\le \eps_*^{(6)}
\le 0.004582638648933954\ldots
<10^{-2}.
\label{eq:k6-bracket-intro}
\end{equation}
The failure is therefore a prospective failure of the selected four-transition witness, not a failure of the unchanged physical map class at six transitions.

The identification consequence is deterministic.  At tolerance $10^{-2}$, the four-transition feasible set contains at least two distinct physical maps with sharply different certified prospective errors.  The accepted map $\Phi_6$ lies in the six-transition feasible set and hence also in the four-transition feasible set, whereas $\Phi_4$ is excluded by the added prospective constraints.  The four finite observations therefore do not uniquely identify a tolerance-feasible physical map.  We use \emph{finite-data non-identification} and \emph{set-membership ambiguity} for this statement.  It does not imply that every possible estimator trained on four transitions must fail, that four transitions are information-theoretically insufficient under every prior or selection rule, or that the effective dynamics is fundamentally unknowable.

The distinction also blocks stronger dynamical readings.  A finite common-map fit is not automatically an autonomous law for arbitrary times, a quantum dynamical semigroup, a CP-divisible process, or a memoryless process.  Likewise, prospective failure of one witness does not prove non-Markovianity, hidden-state necessity, or missing state variables.  Operational approaches to multitime quantum processes make clear that such claims require richer interventions and temporal data \cite{PollockEtAl2018ProcessTensor,PollockEtAl2018Markov,WhiteEtAl2022}.

Our contribution is therefore a controlled finite case rather than a new general identification principle.  We (i) formulate the reviewed calculation as a finite common-channel identification problem; (ii) separate class compatibility, finite-set identification, witness selection, and prospective prediction; (iii) state the reviewed four-transition, prospective-failure, and six-transition certificates in one framework; (iv) use nested feasible sets to make the non-uniqueness deduction explicit; and (v) position this finite BCRD case relative to incomplete quantum process tomography, quantum system identification, finite CPTP interpolation, and classical set-membership identification.  The generic possibility of non-unique model sets under finite or incomplete data is established prior art; the contribution here is the custody-preserved, certified finite quantum instance.

The remainder of the paper is organized as follows.  \Cref{sec:framework} defines the identification problem and feasible-set semantics.  \Cref{sec:physical} specifies the finite record algebra and transition dataset.  \Cref{sec:results} presents the certified K=4 $\rightarrow$ prospective failure $\rightarrow$ K=6 sequence.  \Cref{sec:identification} develops the compatibility-versus-identification interpretation and future feasible-set diagnostics.  \Cref{sec:literature} situates the result within related identification and tomography literature, and \Cref{sec:discussion} states the limitations and nonclaims.
